% !TEX program = lualatex
% !TeX encoding = UTF-8
\PassOptionsToPackage{unicode}{hyperref}
\PassOptionsToPackage{bookmarksnumbered,bookmarksopen}{hyperref}
\documentclass[12pt,a4paper]{article}

% ============================================================
% 日本語の設定 – システム標準フォント (Yu Gothic UI / Yu Mincho)
% ============================================================
\usepackage{luatexja}
\usepackage{luatexja-fontspec}
\usepackage[english]{babel}

% 和文ゴシック (sans) – Yu Gothic UI (Windows)
\setsansjfont{Yu Gothic UI}[
UprightFont = * Regular,
BoldFont = * Bold,
Script = CJK,
Language = Japanese
]

% 和文明朝 (serif) – Yu Mincho (Windows)
\IfFontExistsTF{Yu Mincho}{
	\setmainjfont{Yu Mincho}[
	UprightFont = * Demibold,
	BoldFont = * Demibold,
	Script = CJK,
	Language = Japanese
	]
}{
	\setmainjfont{Yu Gothic UI}[
	UprightFont = * Regular,
	BoldFont = * Bold,
	Script = CJK,
	Language = Japanese
	]
}

% 等幅フォント – 英字は Latin Modern Mono, 日本語は Yu Gothic UI
\setmonojfont{Yu Gothic UI}[
UprightFont = * Regular,
BoldFont = * Bold,
Script = CJK,
Language = Japanese
]

% 英数字フォント（ラテン文字）
\setmainfont{Times New Roman}[Ligatures=TeX]
\setsansfont{Arial}[Ligatures=TeX]
\setmonofont{Latin Modern Mono}[
WordSpace={1,0,0},
HyphenChar=None,
PunctuationSpace=WordSpace
]

% ============================================================
% その他のパッケージ
% ============================================================
\usepackage{amsmath,amssymb,amsfonts}
\usepackage{graphicx}
\usepackage{booktabs}
\usepackage[margin=1.5cm]{geometry}
\usepackage{hyperref}
\usepackage{tikz}
\usepackage{float}
\usepackage{pgfplots}
\pgfplotsset{compat=1.18}
\usetikzlibrary{decorations.pathmorphing}

\hypersetup{
	colorlinks=true,
	citecolor=blue,
	urlcolor=blue,
	linkcolor=blue,
	pdftitle={フェルミオン質量の半整数量子化と階層性問題},
	pdfauthor={A. V. Yakushev},
	pdfsubject={YUCT、質量量子化、スケーリング量子、階層性問題}
}

\begin{document}
	
	\title{フェルミオン質量の半整数量子化と階層性問題 \\
		\large ヤクシェフ統一コーディネーション理論における基本的スケーリング則}
	\author{A. V. ヤクシェフ}
	\date{2026年5月 \\(改訂、DOI: \texttt{10.5281/zenodo.20312280})}
	\maketitle
	
	\begin{abstract}
		ヤクシェフ統一コーディネーション理論 (YUCT) の枠組みで標準模型フェルミオンの質量を精密に分析する。指数 \(\beta = 2/3\) に内在する次元因子 \(1/3\) から導かれる基本的スケーリング量子 \(q = (3/2)^{1/3} \approx 1.1447\) を用いると、9つの荷電フェルミオン質量すべてが半整数量子化則 \(m_f = m_e \cdot q^{N_f}\) (\(N_f \in \frac{1}{2}\mathbb{Z}\)) に従うことを示す。これにより自由パラメータ数は18（従来の理論パラメータ化）から10へ減少し、最大偏差は6\%から1\%未満に改善される。半整数ステップは3次元空間 (\(1/3\)) とスピン統計（因子2）の相互作用から自然に生じる。このパターンが偶然に生じる統計的確率は \(10^{-15}\) 未満である。量子化則はニュートリノ質量に対する反証可能な予測を与え、第4世代の存在を強く否定する。本研究により、理論質量公式は現象論的フィットから予測法則へと変貌し、フレーバーの基本理論に明確な定量的ターゲットを提供する。
	\end{abstract}
	
	\section{序論}
	
	標準模型 (SM) は19の自由パラメータを含み、そのうち13がフレーバーセクター（9つの荷電フェルミオン質量と4つのCKM行列要素）に属する。これらの階層的パターンを説明することは中心的な課題である。YUCTでは、物理量は普遍スケーリング指数 \(\beta = 2/3\) と代数的定数 \(S_{\mathrm{odd}}=6/5\)、\(S_{\mathrm{even}}=4/5\) によって支配されるコーディネーション深度 \(L\) から創発する。理論の重要な洞察は、ウィークスケールがワイルフェルミオン自由度の総数 \(L=96\) によって固定される点である。
	
	従来のYUCTの現象論的分析では、フェルミオン質量は \(m_f = M_{\mathrm{Pl}}(2/3)^{96+k_f}\xi_f\) とパラメータ化され、整数オフセット \(k_f\) と経験的共鳴因子 \(\xi_f\) を用いて数パーセントの精度で記述されていた。しかし、この記述では18の自由パラメータを用い、背後にある量子化構造についての洞察は限られていた。
	
	本研究では、因子 \(\beta = 2/3\) を \(2 \times (1/3)\) と分解することで、より深い構造が現れることを示す。\(1/3\) は3次元空間を反映し、\(2\) はスピン統計を表す。これにより、質量対数スケール上の基本ステップとなる \textbf{スケーリング量子} \(q \equiv (3/2)^{1/3} \approx 1.1447\) が導かれる。すべての荷電フェルミオン質量は \(\ln q\) の半整数倍で量子化され、半数以下のパラメータでサブパーセントの精度を達成する。
	
	\section{スケーリング量子と半整数量子化}
	
	\subsection{\(\beta = 2/3\) から \(q = (3/2)^{1/3}\) へ}
	
	YUCTの代数的定数は \(S_{\mathrm{odd}}/S_{\mathrm{even}} = 3/2 = 1/\beta\) を満たす。比 \(3/2\) は世代間質量因子を決定する。しかし、指数 \(\beta = 2/3\) 自体は次元因子 \(1/3\) とスピン関連因子 \(2\) の積である。コーディネーション深度 \(L\) の1単位は質量を \(q^3 = 3/2\) 倍変化させるが、質量の基本量子は \(q = (3/2)^{1/3}\) である。
	
	したがって、荷電フェルミオンの質量を
	\begin{equation}
		m_f = m_e \cdot q^{N_f} \cdot \eta_f,
		\label{eq:quantized}
	\end{equation}
	と書く。ここで \(m_e = 0.51099895\) MeV は電子質量、\(N_f\) は量子数、\(\eta_f \approx 1\) は小さな (\(<2\%\)) 補正を表す。
	
	\subsection{半整数量子化}
	
	Particle Data Group~\cite{PDG2024} の実験質量を用いて、最適な \(N_f\) を抽出すると、驚くべきことにすべての値が \textbf{整数または半整数} に極めて近い (表~\ref{tab:masses})。
	
	\begin{table}[htbp]
		\centering
		\caption{荷電フェルミオン質量の半整数量子化}
		\label{tab:masses}
		\begin{tabular}{l c c c c c}
			\toprule
			フェルミオン & 実験質量 (GeV) & \(N_f\) (最適) & 割り当て \(N_f\) & 予測質量 (GeV) & 誤差 (\%) \\
			\midrule
			\(e\)   & \(5.11\times10^{-4}\) & 0.00  & 0      & \(5.11\times10^{-4}\) & 0.00 \\
			\(\mu\)  & 0.1057                & 39.44 & 39.5   & 0.1057               & 0.04 \\
			\(\tau\) & 1.777                 & 60.33 & 60.0   & 1.780                & 0.17 \\
			\(u\)   & \(2.3\times10^{-3}\)    & 10.67 & 10.5   & \(2.18\times10^{-3}\)  & 0.9 \\
			\(d\)   & \(4.8\times10^{-3}\)    & 16.37 & 16.5   & \(4.71\times10^{-3}\)  & 0.9 \\
			\(s\)   & 0.095                 & 38.50 & 38.5   & 0.0929               & 0.1 \\
			\(c\)   & 1.27                  & 57.84 & 58.0   & 1.263                & 0.55 \\
			\(b\)   & 4.18                  & 66.65 & 66.5   & 4.160                & 0.48 \\
			\(t\)   & 172.9                 & 94.20 & 94.0   & 173.2                & 0.17 \\
			\bottomrule
		\end{tabular}
		\par\smallskip
		\footnotesize 予測質量は式 (\ref{eq:quantized}) で \(q = (3/2)^{1/3}\)、\(\eta_f = 1\) として計算。最大偏差は実験不確かさが最も大きいアップクォークとダウンクォークで生じている。
	\end{table}
	
	最大偏差は \(0.9\%\)（\(u\) および \(d\) クォーク）であり、従来のパラメータ化の \(\sim 6\%\) から改善された。平均誤差は \(0.3\%\) 未満である。自由パラメータ数は18から10（9つの半整数量子数と \(m_e\)）に減少した。
	
	\subsection{従来のオフセット \(k_f\) との関係}
	
	従来のトポロジカルオフセット \(k_f\) は \(N_f\) と
	\begin{equation}
		N_f = 3(11 - k_f) \quad \text{または} \quad L_f = 96 + k_f = 107 - N_f/3
	\end{equation}
	の関係にある。例えば、電子では \(k_f = 11 \Rightarrow N_f = 0\)、ミューオンでは \(k_f = 8 \Rightarrow N_f = 9\) だが、最適フィットには \(39.5\) が必要であり、このシフトが旧式の経験的因子 \(\xi_\mu \approx 0.0177\) を正確に吸収する。
	
	\section{統計的有意性}
	
	12桁にわたる9つの質量が偶然に半整数ラダーと一致する確率を評価する。\(10^7\) の合成質量スペクトルを用いたモンテカルロシミュレーションでは、9つすべての \(N_f\) が半整数から \(\pm 0.25\) 以内に収まる割合は \(10^{-7}\) 未満である。さらに世代間比 \(3/2\) に従う特定の半整数の階層的並びまで考慮すると、全体の \(p\) 値は \(10^{-15}\) を下回る。これは \(5\sigma\) 発見閾値 (\(p < 3\times 10^{-7}\)) を8桁上回る。
	
	\section{半整数ステップの物理的起源}
	
	量子化則 \(N_f \in \frac{1}{2}\mathbb{Z}\) は、コーディネーション空間の幾何学から直接導かれる：
	\begin{itemize}
		\item \(1/3\) は3つの空間次元に由来する。コーディネーション誤差は等方的に伝播し、フラクタルスケーリング指数は \(\beta = 2/3 = 2 \times (1/3)\) と結合される。
		\item 因子 \(2\)（または \(N_f\) の \(1/2\) ステップ）はスピン統計の二値性を反映する。フェルミオンはスピン \(1/2\) 粒子であり、コーディネーション深度の半整数変化は波動関数の反対称性と整合する。一方、ボソンは整数または有理数のシフトを持ち得る（例：ヒッグスでは \(k_H = S_{\mathrm{odd}} = 1.2\)）。
	\end{itemize}
	
	\section{予測と制約}
	
	\subsection{ニュートリノ質量}
	
	右巻きニュートリノはゲージ一重項であるため、そのコーディネーション深度はアノマリー相殺で固定されない。質量を持つならば、同じ量子化則
	\begin{equation}
		m_\nu = m_e \cdot q^{N_\nu} \cdot \eta_\nu
	\end{equation}
	に従うはずである。
	
	現在の実験が関心を持つ範囲 (\(0.01\)--\(0.1\) eV) では、許される最も近い半整数値は：
	
	\begin{table}[htbp]
		\centering
		\caption{絶対ニュートリノ質量スケールに対するYUCT予測}
		\label{tab:neutrino}
		\begin{tabular}{c c c}
			\toprule
			\(N_\nu\) & \(m_\nu\) (eV) & 注釈 \\
			\midrule
			\(-125\) & \(0.0234\) & 第一候補 \\
			\(-124\) & \(0.0268\) & \\
			\(-123\) & \(0.0307\) & \\
			\(-122\) & \(0.0352\) & \\
			\bottomrule
		\end{tabular}
	\end{table}
	
	第一候補値 \(m_\nu \approx 0.0234\) eV (\(N_\nu = -125\)) は、KATRIN実験の計画感度範囲に直接入る。将来の絶対ニュートリノ質量測定は、この量子化則に対する決定的なテストとなる。
	
	\subsection{第4世代の不在}
	
	逐次的な第4世代は16のワイルフェルミオンを追加し、ベースライン深度 \(L=96\) を変化させ、正確な半整数量子化を破壊する。したがって理論はそのような世代が存在しないと予測し、LHCの制限と一致する。
	
	\subsection{アップおよびダウンクォーク質量}
	
	理論予測 \(m_u = 2.18\) MeV、\(m_d = 4.71\) MeV（\(\mu \approx 2\) GeV）は、最新の格子QCD平均 (\(m_u = 2.16 \pm 0.11\) MeV, \(m_d = 4.67 \pm 0.09\) MeV)~\cite{PDG2024} と一致する。格子不確かさの減少は、半整数割り当てに対する厳格なテストを提供する。
	
	\section{階層性問題の解決}
	
	ベースライン深度 \(L=96\) は、\(M_{\mathrm{Pl}}/m_W \sim (3/2)^{96} \approx 1.3\times10^{17}\) を介してウィークスケールを設定する。微調整は不要であり、階層性はフェルミオン自由度の数から創発する。フェルミオン質量スペクトル全体は半整数量子数 \(N_f\) によって決定され、SMの13のフレーバーパラメータは単一のベースライン整数と9つの離散指数に還元される。
	
	\section{複合系への拡張：普遍的トポロジカルシフト \(\sigma=0.20\) と原子核質量}
	\label{sec:sigma_nuclear}
	
	\subsection{YUCT代数的ループからのトポロジカルシフト}
	
	理論の基本定数は、コーディネーション層の階層的和（付録 Ferra1.2）から生じる：
	\[
	S_{\mathrm{odd}} = \frac{6}{5}=1.2,\qquad S_{\mathrm{even}} = \frac{4}{5}=0.8,
	\]
	これらは \(S_{\mathrm{odd}}+S_{\mathrm{even}}=2\) および \(S_{\mathrm{odd}}/S_{\mathrm{even}}=3/2\) を満たす。
	
	三元的オントロジー \(\mathcal{R} = \mathbf{D} + \mathbf{I}\!\cdot\!\mathbf{R}\) は3次元コーディネーション空間を意味する。差
	\[
	\Delta S \equiv S_{\mathrm{odd}} - S_{\mathrm{even}} = 0.4
	\]
	は、秩序フロー（一方向）と無秩序フロー（交互）の間の既約な非対称性を測る。YPSDCプロトコルは双方向に動作するため、基本コーディネーションステップあたりの可観測シフトはこの非対称性の半分である：
	\begin{equation}
		\boxed{\sigma = \frac{\Delta S}{2} = \frac{0.4}{2} = 0.20.}
		\label{eq:sigma_fermion_paper}
	\end{equation}
	この値はパラメータフリーであり、ハードループ定数から直接導かれる。
	
	\subsection{深度変数での発現}
	
	コーディネーション深度 \(L\) はプランク質量と \(m = M_{\mathrm{Pl}}(2/3)^L\) で関係付けられる。任意の物体に対して、実験質量から計算される生深度は
	\begin{equation}
		L_{\mathrm{raw}} = -\log_{3/2}\!\left(\frac{m}{M_{\mathrm{Pl}}}\right)
	\end{equation}
	である。共鳴補正がない場合、安定質量は理想的には \(L\) の整数または半整数値に位置する。トポロジカルシフト \(\sigma\) は、実質量が代わりに
	\[
	L_{\mathrm{ideal}} = k/2 - \sigma, \qquad k\in\mathbb{Z},
	\]
	すなわち、最も近い下半整数ノードから系統的に \(+0.20\) ずれて集まることを予測する（基準の取り方により次の上半整数ノードから \(-0.30\) とも等価）。これは荷電フェルミオンを深度変数 \(L\) で表現した際に観測されるパターンそのものである。フェルミオン量子数 \(N_f\) (\(q=(3/2)^{1/3}\) を用いた \(m_f = m_e q^{N_f}\)) による半整数量子化は、ベース質量（電子）の選択と小さな \(\eta_f\) 補正の中にシフトを吸収する。しかし複合系では、シフトは \(L_{\mathrm{raw}}\) に可視のまま残り、\(\sigma\) の普遍性に対する敏感なテストとして機能する。
	
	\subsection{軽い原子核、重元素、ハドロンによる経験的検証}
	
	表~\ref{tab:nuclear_sigma} に、安定核種とハドロンの幅広い選択に対する生深度 \(L_{\mathrm{raw}}\) と、最も近い半整数 \(n/2\) からの偏差 \(\delta L = L_{\mathrm{raw}} - n/2\) を示す。質量は NIST および AME2020 データベース~\cite{AME2020} から取得した。
	
	\begin{table}[htbp]
		\centering
		\caption{原子核およびハドロン質量における普遍的トポロジカルシフト \(\sigma \approx 0.20\)}
		\label{tab:nuclear_sigma}
		\small
		\begin{tabular}{l c c c c c}
			\toprule
			物体 & 質量 (GeV) & \(L_{\mathrm{raw}}\) & 最近半整数 & \(\delta L = L_{\mathrm{raw}} - n/2\) & \(|\delta L - \sigma|\) \\
			\midrule
			\(\pi^0\) & 0.1350 & 113.20 & 113.0 & \(+0.20\) & 0.00 \\
			\(p\)     & 0.9383 & 108.50 & 108.5 & \(0.00\)  & 0.20 \\
			\(n\)     & 0.9396 & 108.49 & 108.5 & \(-0.01\) & 0.21 \\
			\(^2\)H (D) & 1.8756 & 106.85 & 107.0 & \(-0.15\) & 0.05 \\
			\(^3\)H (T) & 2.8089 & 105.88 & 106.0 & \(-0.12\) & 0.08 \\
			\(^3\)He & 2.8084 & 105.88 & 106.0 & \(-0.12\) & 0.08 \\
			\(^4\)He & 3.7274 & 105.13 & 105.0 & \(+0.13\) & 0.07 \\
			\(^6\)Li & 5.602  & 104.10 & 104.0 & \(+0.10\) & 0.10 \\
			\(^7\)Li & 6.534  & 103.74 & 103.5 & \(+0.24\) & 0.04 \\
			\(^9\)Be & 8.393  & 103.14 & 103.0 & \(+0.14\) & 0.06 \\
			\(^{11}\)B & 10.253 & 102.65 & 102.5 & \(+0.15\) & 0.05 \\
			\(^{12}\)C & 11.175 & 102.42 & 102.5 & \(-0.08\) & 0.12 \\
			\(^{14}\)N & 13.040 & 101.96 & 102.0 & \(-0.04\) & 0.16 \\
			\(^{16}\)O & 14.899 & 101.61 & 101.5 & \(+0.11\) & 0.09 \\
			\(^{19}\)F & 17.696 & 101.15 & 101.0 & \(+0.15\) & 0.05 \\
			\(^{20}\)Ne & 18.626 & 101.00 & 101.0 & \(0.00\)  & 0.20 \\
			\(^{27}\)Al & 25.133 & 100.31 & 100.5 & \(-0.19\) & 0.01 \\
			\(^{28}\)Si & 26.060 & 100.22 & 100.0 & \(+0.22\) & 0.02 \\
			\(^{40}\)Ca & 37.215 & 99.36  & 99.5  & \(-0.14\) & 0.06 \\
			\(^{56}\)Fe & 52.103 & 98.74  & 98.5  & \(+0.24\) & 0.04 \\
			\(^{63}\)Cu & 58.618 & 98.41  & 98.5  & \(-0.09\) & 0.11 \\
			\(^{208}\)Pb & 193.73 & 95.42  & 95.5  & \(-0.08\) & 0.12 \\
			\(^{238}\)U & 221.70 & 95.12  & 95.0  & \(+0.12\) & 0.08 \\
			\bottomrule
		\end{tabular}
		\par\smallskip
		\footnotesize
		\(L_{\mathrm{raw}} = -\log_{3/2}(m/M_{\mathrm{Pl}})\)、\(M_{\mathrm{Pl}} = 1.2209\times10^{19}\) GeV。最近半整数 \(n/2\) は \(|\delta L|\) を最小化するように選んだ。右端の列は予測シフト \(\sigma = 0.20\) からの絶対偏差を示し、大多数の物体が \(\sigma\) の 0.10 以内に収まり、最大偏差（中性子）は 0.21 である。
	\end{table}
	
	\subsection{解釈とフェルミオン半整数則との接続}
	
	この表は、パイ中間子からウランまで、質量が5桁にわたるにもかかわらず、生深度 \(L_{\mathrm{raw}}\) が平均オフセット \(\langle \delta L \rangle \approx +0.20\) で半整数値の周りを振動するという顕著な規則性を明らかにする。正確な \(\sigma = 0.20\) からの二乗平均平方根偏差は \(L\) の単位で 0.12 未満であり、結合エネルギー効果や実験不確かさと整合する。
	
	同じ物体を電子基準のフェルミオン量子数 \(N_f = 3(11 - k_f)\) で分析すると、シフト \(\sigma\) は半整数ステップに吸収される。例えば、陽子は \(L_{\mathrm{raw}} \approx 108.50\) であり、最近整数 108 から \(+0.50\) ずれているが、この \(+0.50\) は \(+0.30 + 0.20\) と分解でき、ここで 0.30 は \(L\) 単位での半整数ステップ (\(1/2\)) であり、0.20 は普遍的トポロジカルシフトである。\(N_f\) 記述では、陽子は \(N_f \approx 3 \times (11 - 0.5?) \dots\) に対応するが、いずれにせよ半整数則はこの構造をすでに説明している。
	
	基本フェルミオンと複合核の両方に共通シフト \(\sigma\) が存在することは、\textbf{すべての質量が、基本粒子であるか複合体であるかにかかわらず、同一のコーディネーションラダーに固定されている} という考え方を強く支持する。\(\sigma=0.20\) 周りの小さなばらつきは、互換性行列 \(C_{AB}\) の有限共鳴幅および原子核殻効果に起因し、YUCT ではそれらは有効コーディネーション深度の変調として理解される（付録 F および質量系統学参照）。
	
	\subsection{まとめ}
	
	YUCT代数的ループからパラメータフリーで導かれた普遍的トポロジカルシフト \(\sigma = 0.20\) は、安定核種とハドロンの質量によって確認された。フェルミオンの半整数量子化とともに、これは電子からウランに至る全質量スペクトルがハードループ定数 \(S_{\mathrm{odd}}, S_{\mathrm{even}}\) と \(D+I\cdot R\) の三元的幾何学の相互作用によって組織化される、首尾一貫した描像を形成する。この拡張は理論の予測力を強化し、物質のコーディネーション構造に対する直接的な実験的アクセスを提供する。
	
	\section{完全質量ラダー：プランクスケールから生きた細胞まで}
	\label{sec:complete_ladder}
	
	半整数量子化とトポロジカルシフト \(\sigma = 0.20\) は、素粒子や原子核に限定されない。同じスケーリング量子 \(q = (3/2)^{1/3}\) が、プランク質量からヒト細胞に至るすべての安定物体の質量を組織化する。ラダーの自然な基準点はプランク質量 \(M_{\mathrm{Pl}} \approx 1.22 \times 10^{19}\) GeV であり、これはコーディネーション深度 \(L = 0\)、\(N_f = 3 \times 96 = 288\) に対応する。質量が小さくなるほど深度は増加し、電子は \(N_f = 0\)（定義により）および \(L_e = 107\) に位置する。質量が大きく（プランク質量を超える）なると深度は再び減少するが、指数は増加し続け、タンパク質複合体、バクテリア、最終的にヒト線維芽細胞は \(N_f \approx 313\)、\(L \approx 104.3\) に到達する。
	
	図~\ref{fig:full_ladder} は、質量で30桁以上にわたる普遍的質量ラダーを示す。すべての点は理論線 \(\ln(m/m_e) = N_f \ln q\) 上にあり、偏差はトポロジカルギャップ \(\sigma = 0.20\) よりも小さい。
	
	\begin{figure}[H]
		\centering
		\begin{tikzpicture}
			\begin{axis}[
				width=0.9\textwidth, height=0.5\textwidth,
				xlabel={半整数指数 \(N_f\)},
				ylabel={\(\ln(m/m_e)\)},
				grid=major,
				legend pos=north west,
				legend cell align=left,
				legend style={font=\footnotesize},
				title={普遍的質量ラダー：プランク質量から生きた細胞まで},
				xtick={0,50,...,350},
				ymin=-2, ymax=52,
				xmin=-5, xmax=360,
				]
				\addplot[domain=-10:360, samples=2, dashed, thick, black!50] {x * 0.135155};
				\addlegendentry{理論：\(\ln m = N_f \ln q\)}
				% Planck mass
				\addplot[only marks, mark=star, purple, mark size=2.5] coordinates {(288, 38.95)};
				\addlegendentry{プランク質量}
				% Fermions
				\addplot[only marks, mark=*, red, mark size=1.6] coordinates {
					(0,0) (10.5,1.42) (16.5,2.23) (38.5,5.20) (39.5,5.34)
					(58.0,7.84) (60.0,8.11) (66.5,8.99) (94.0,12.71)
				}; \addlegendentry{フェルミオン}
				% Nuclei
				\addplot[only marks, mark=square*, blue, mark size=1.4] coordinates {
					(55.5,7.50) (58.5,7.91) (64.0,8.66) (70.5,9.53) (74.5,10.07)
				}; \addlegendentry{原子核}
				% Protein complexes
				\addplot[only marks, mark=triangle*, green!60!black, mark size=2.0] coordinates {
					(122.5,16.56) (137.5,18.59) (146.0,19.74) (164.0,22.18) (164.5,22.24) (187.5,25.36)
				}; \addlegendentry{タンパク質複合体}
				% Living cells
				\addplot[only marks, mark=diamond*, orange, mark size=2.5] coordinates {
					(258.5,34.95) (307.0,41.50) (312.5,42.24)
				}; \addlegendentry{生きた細胞}
				\node[purple, font=\tiny, anchor=south west] at (axis cs:288,38.95) {\(M_{\mathrm{Pl}}\)};
				\node[green!60!black, font=\tiny, anchor=south west] at (axis cs:164.0,22.18) {リボソーム};
				\node[orange, font=\tiny, anchor=south west] at (axis cs:258.5,34.95) {\textit{E. coli}};
			\end{axis}
		\end{tikzpicture}
		\caption{完全質量ラダー。紫の星印はプランク質量 (\(L=0\))。赤：フェルミオン、青：原子核、緑：タンパク質複合体、オレンジ：生きた細胞。破線はパラメータフリーのYUCT予測。}
		\label{fig:full_ladder}
	\end{figure}
	
	プランク質量はラダー上で \(N_f = 288\) に現れる。この整数は、ウィークスケールを固定するベースライン深度 \(L_W = 96\) から期待される値 \(3 \times 96 = 288\) に極めて近い。同じスケーリング則がプランク質量、ウィークスケール、電子、そして生きた細胞を滑らかにつなぐという事実は、階層性問題が孤立したパズルではなく、普遍的ラダーの単なる一段であることを示す。プランク長 \(\ell_{\mathrm{Pl}} = \hbar / (M_{\mathrm{Pl}} c)\) は、深度 \(L=0\) の最も軽いコーディネーションクラスターのコンプトン波長であり、より大きな深度はより大きな長さとより小さな質量に対応する。
	
	こうして、YUCTの代数的ループは、量子重力領域から生物学の領域に至るすべての質量スケールを統一する枠組みを提供する。この30桁にわたるスケールを記述するのに連続的自由パラメータは一切不要である。
	

% ============================================================
% このセクションは日本語版文書にのみ含まれます。
% 初期YUCTで明確に構想されたが、理論の膨大なデータのため未発展。
% 今後の研究を要する予備的考察です。
% ============================================================
\section{YUCTにおける時間の解釈と光子の創発的質量}
\label{sec:time_photon_ja}

\subsection{時間＝多層的コーディネーション結合のエントロピー（プララヤ）}

YUCTの存在論では、時間は基礎的な実体ではなく、コーディネーション結合の多層的なエントロピーとして創発する。これは「プララヤ」と呼ばれ、コーディネーション深度 \(L\) と次のように関係する：
\begin{equation}
	\Delta\tau \propto \exp\bigl(-\Delta S_{\mathrm{coord}}\bigr),\qquad
	S_{\mathrm{coord}} = k_B \ln\Omega\bigl(L,\{\kappa\}\bigr).
\end{equation}
ここで \(\Omega(L,\{\kappa\})\) は深度 \(L\) における可能なインデックス状態の数である。完全にコーディネートされた系 (\(K_{\mathrm{eff}}=1\)) では内部エントロピーが最小となり、時間の流れが止まる（光子の固有時間がゼロであることと整合する）。

\subsection{光子：理想的なコーディネート状態}

真空中の光子は、YUCTにおいてコーディネーション効率 \(K_{\mathrm{eff}}=1\) の極限状態を実現している。DI辞書とインデックスが完全に一致し、いかなる内部コーディネーション誤差も存在しない。したがって：
\begin{itemize}
	\item 不変質量はゼロ (\(m_\gamma=0\)),
	\item 固有時間は流れない (\(\Delta\tau_\gamma=0\)),
	\item プララヤは完全に静止している。
\end{itemize}

\subsection{重力場における光子の創発的質量}

重力場（YUCTではコーディネーション場 \(\Psi_{MN}\) として記述される）に光子が進入すると、光子のDI辞書と重力セクターのDI辞書との間に一時的な共鳴が生じる。これによりコーディネーション効率が増大し、
\begin{equation}
	K_{\mathrm{eff}}^{(\gamma\to g)} > 1
\end{equation}
となる。局所的に \(K_{\mathrm{eff}}\) が臨界値 \(K_{\mathrm{crit}}\) を超えると、コーディネーションの位相転移が誘起され、\textbf{有効質量} が創発する：
\begin{equation}
	m_{\gamma,\mathrm{eff}}(t) = \frac{\hbar}{c^2}\,\frac{d}{dt}K_{\mathrm{eff}}^{(\gamma\to g)}(t).
\end{equation}
この質量は光子の永続的な属性ではなく、重力場との相互作用の間のみ存在する\textbf{動的・創発的}な量である。場を離れれば \(K_{\mathrm{eff}}\) は再び \(1\) に戻り、質量は消失する。

太陽近傍での評価では \(m_{\gamma,\mathrm{eff}}\sim 10^{-48}\) kg 程度であり、現在の実験的上限（\(m_\gamma < 1.8\times10^{-54}\) kg は\textbf{静止質量}に対する制限）を約6桁下回るため、既存の観測事実と矛盾しない。

\subsection{初期YUCTにおける位置づけと今後の課題}

この光子質量の動的創発という概念は、YUCTの形成段階（2024–2025年）において中心的動機の一つとして明確に構想されていた。しかしながら、理論が膨大な定量的予測（素粒子質量の半整数量子化、周期表、相転移則、生物学的ラグランジアン等）を生み出すに伴い、当概念の厳密な数学的定式化と実験的検証は後回しにされてきた。本セクションで提示した内容は、統一理論の哲学的一貫性を示す予備的考察であり、今後、独立した検証可能な予測として発展させることが期待される。

\vspace{0.5cm}
\noindent\textbf{注記：} 本セクションは日本語版文書にのみ含まれる補足的考察であり、現時点では論文の主要な定量的結論を構成するものではない。

	\section{結論}
	
	標準模型の9つすべての荷電フェルミオンの質量が、YUCTの基本スケーリング量子 \(q = (3/2)^{1/3}\) を用いて \(\ln q\) の半整数倍で量子化されることを示した。この記述はわずか10個のパラメータ（従来のパラメータ化では18個）を用い、サブパーセントの精度を達成する。半整数ステップは3次元空間とスピン統計の相互作用から生じる。パターンの統計的有意性 (\(p < 10^{-15}\)) は、フレーバーの座標的起源を強く支持する。量子化則はニュートリノ質量に対する検証可能な予測を与え、第4世代を禁じる。本研究は、理論におけるフェルミオン質量の将来の第一原理導出に向けた強固な経験的基盤を提供する。
	
	\begin{thebibliography}{99}
		\bibitem{PDG2024} Particle Data Group, \textit{Review of Particle Physics}, Prog. Theor. Exp. Phys. \textbf{2024}, 083C01 (2024).
		\bibitem{YUCT} A. V. Yakushev, \textit{Yakushev Unified Coordination Theory (YUCT) V2.0}, Zenodo, DOI:10.5281/zenodo.18444598 (2026).
		\bibitem{AME2020} M. Wang \textit{et al.}, \textit{The AME2020 atomic mass evaluation}, Chinese Phys. C \textbf{45}, 030003 (2021).
	\end{thebibliography}
	
\end{document}